\section*{1. Inverses and the inverse algorithm}

\notetarget{note:02:inverse-definition}An \textbf{inverse} of \(A\in\R^{n\times n}\) is an \(n\times n\) matrix
\(B\) satisfying both \(AB=I_n\) and \(BA=I_n\). A matrix with an
inverse is \textbf{invertible}. Its inverse is unique: if \(B,C\) are
inverses, associativity gives
\[
B=BI=B(AC)=(BA)C=IC=C.
\]
Write this unique matrix as \(A^{-1}\). Since
\(AA^{-1}=A^{-1}A=I_n\), \(A\) is the inverse of \(A^{-1}\):
\((A^{-1})^{-1}=A\).

\begin{proposition}[Product and transpose inverses]\label{prop:02:inverse-laws}\notetarget{note:02:inverse-product}
If \(A,B\) are invertible square matrices of the same size, then
\[
(AB)^{-1}=B^{-1}A^{-1},\qquad (A^\top)^{-1}=(A^{-1})^\top.
\]
\end{proposition}
\begin{proof}
The proposed inverse satisfies both defining identities:
\[
(AB)(B^{-1}A^{-1})=A(BB^{-1})A^{-1}=I,\qquad
(B^{-1}A^{-1})(AB)=B^{-1}(A^{-1}A)B=I.
\]
\notetarget{note:02:inverse-transpose}Transposing \(AA^{-1}=I\) and \(A^{-1}A=I\), using
Unit 1, Proposition~\ref{one-prop:01:transpose}, gives
\[
(A^{-1})^\top A^\top=I,\qquad A^\top(A^{-1})^\top=I.
\]
Thus the proposed matrices are inverses by definition.
\end{proof}
Repeated application of the product formula gives
\((A_k\cdots A_1)^{-1}=A_1^{-1}\cdots A_k^{-1}\): append one factor
and apply the two-factor formula, starting with a single factor.
\notetarget{note:02:cancellation}If \(A\) is invertible, then \(AX=AY\) implies \(X=Y\), by multiplying
on the left by \(A^{-1}\). Similarly \(XA=YA\) permits multiplication
on the right by \(A^{-1}\). The side of multiplication matters.

\subsection*{Recording a row operation}
\notetarget{note:02:elementary-matrices}An \textbf{elementary matrix} is obtained by applying one elementary
row operation to \(I_m\). It records that operation by left multiplication.

\begin{proposition}[Elementary matrices]\label{prop:02:elementary}
If \(E\) is obtained by a row operation on \(I_m\), then \(EA\) is obtained
by the same operation on every \(m\)-row matrix \(A\).
Each \(E\) is invertible, with inverse given by the reverse operation.
\end{proposition}
\begin{proof}
Row \(i\) of \(EA\) is \(\sum_k e_{ik}\,\operatorname{row}_k(A)\):
its entry in column \(j\) is the product formula \(\sum_k e_{ik}a_{kj}\).
Thus the three row operations on \(I_m\) give the following rows of \(EA\):
\begin{itemize}
\item Swap \(r,s\): rows \(r,s\) become
\(\operatorname{row}_s(A),\operatorname{row}_r(A)\).
\item Scale row \(r\) by \(c\ne0\): row \(r\) becomes
\(c\operatorname{row}_r(A)\).
\item Add \(cR_s\) to \(R_r\), \(r\ne s\): row \(r\) becomes
\(\operatorname{row}_r(A)+c\operatorname{row}_s(A)\).
\end{itemize}
All other rows are unchanged.

\notetarget{note:02:elementary-reverse}The reverse operations are the same swap, scaling by \(1/c\), and
adding \(-c\) times the unchanged source row. Let \(F\) record the
reverse operation. Applying either order to \(I_m\) gives
\(FE=EF=I_m\), so \(F=E^{-1}\).
\end{proof}
\notetarget{note:02:elementary-example}For example, adding \(c\) times row 1 to row 3 of a three-row matrix
is recorded by
\[
E=\begin{pmatrix}1&0&0\\0&1&0\\c&0&1\end{pmatrix},\qquad
E^{-1}=\begin{pmatrix}1&0&0\\0&1&0\\-c&0&1\end{pmatrix}.
\]
The change \(c\mapsto-c\) records the reverse row operation.

\notetarget{note:02:row-record}After row operations \(E_1,\ldots,E_k\), the result is
\[
H=E_k\cdots E_1A=EA,\qquad E^{-1}=E_1^{-1}\cdots E_k^{-1}.
\]
Thus \(E\) is invertible. A row operation is a reversible change in the
equations; the product \(E\) records all the changes in their actual order.

\begin{theorem}[Square invertibility criterion]\label{thm:02:inverse}\notetarget{note:02:square-criterion}
For \(A\in\R^{n\times n}\), the following are equivalent:
\begin{enumerate}
\item \(A\) is invertible.
\item \(A\) is row equivalent to \(I_n\).
\item \(Ax=0\) has only the solution \(x=0\).
\item For every \(b\in\R^n\), \(Ax=b\) has exactly one solution.
\end{enumerate}
\end{theorem}
\begin{proof}
\begin{itemize}
\item (a) \(\Longrightarrow\) (c): \(Ax=0\) gives \(x=A^{-1}Ax=0\).
\item \notetarget{note:02:square-pivots}(c) \(\Longrightarrow\) (b): by Unit 1, Theorems~\ref{one-thm:01:echelon}
and~\ref{one-thm:01:consistency}, elimination can leave no free variable.
There are therefore \(n\) pivots. Their strictly increasing column
indices must be \(1,\ldots,n\), so the echelon matrix is upper triangular:
all entries below the diagonal are zero. Its diagonal entries are nonzero.
Divide each row by its diagonal entry.
\notetarget{note:02:square-backward}The last row is now \(e_n^\top\); subtract suitable multiples of it from
all earlier rows to clear the last column. The next-to-last row is then
\(e_{n-1}^\top\); clear its column above the diagonal. Continuing upwards
produces \(I_n\). Each operation is permitted, proving (b).

\item \notetarget{note:02:square-cycle}(b) \(\Longrightarrow\) (a): record the reduction as \(EA=I_n\), where \(E\) is a
product of elementary matrices. By Proposition~\ref{prop:02:elementary}
and the inverse product law, \(E\) is invertible. Multiplying by
\(E^{-1}\) gives \(A=E^{-1}\). Hence \(A\) has inverse \(E\).

\item (a) \(\Longrightarrow\) (d): \(A(A^{-1}b)=b\) proves existence;
\(Ax=b\) forces \(x=A^{-1}b\), proving uniqueness.
\item (d) \(\Longrightarrow\) (c): take \(b=0\); its unique solution is zero.
\end{itemize}
\end{proof}
\notetarget{note:02:elementary-factorization}In particular an invertible matrix is a product of elementary matrices:
if \(E_k\cdots E_1A=I\), then
\(A=E_1^{-1}\cdots E_k^{-1}\), and each inverse is elementary.
Conversely any such product is invertible by the product law.

\begin{proposition}[One-sided inverses for square matrices]\label{prop:02:one-sided}\notetarget{note:02:one-sided}
For square \(A,B\) of the same size, either \(BA=I\) or \(AB=I\)
implies that \(B=A^{-1}\).
\end{proposition}
\begin{proof}
If \(BA=I\) and \(Ax=0\), then \(x=BAx=0\).
Theorem~\ref{thm:02:inverse} makes \(A\) invertible, and
\(B=B(AA^{-1})=(BA)A^{-1}=A^{-1}\).
If \(AB=I\), interchange \(A,B\) in the preceding argument. It gives
\(A=B^{-1}\), so \(B=(B^{-1})^{-1}=A^{-1}\).
\end{proof}
\notetarget{note:02:rectangular}For rectangular matrices one identity need not imply the other:
\[
\begin{pmatrix}1&0\end{pmatrix}\begin{pmatrix}1\\0\end{pmatrix}=(1),
\qquad
\begin{pmatrix}1\\0\end{pmatrix}\begin{pmatrix}1&0\end{pmatrix}
=\begin{pmatrix}1&0\\0&0\end{pmatrix}\ne I_2.
\]
These are the standard-coordinate inclusion and selection matrices.

\subsection*{Computing an inverse}
\notetarget{note:02:inverse-algorithm}Apply the same row operations to both halves of \([A\mid I_n]\).
If the left half becomes \(I_n\), the result is
\[
[A\mid I_n]\longmapsto[EA\mid E]=[I_n\mid A^{-1}].
\]
The justification is exactly the implication (b) to (a) above.
If elimination of \(A\) leaves a free variable, then \(N(A)\ne\{0\}\)
by back substitution, so \(A\) is not invertible.

\begin{example}[An inverse by row operations]\label{example:02:inverse}\notetarget{note:02:inverse-example}\notetarget{note:02:inverse-elimination}
For \(A=\begin{pmatrix}0&1&1\\1&1&1\\3&4&3\end{pmatrix}\),
\begin{align*}
\left[\begin{array}{rrr|rrr}0&1&1&1&0&0\\1&1&1&0&1&0\\3&4&3&0&0&1\end{array}\right]
&\xrightarrow{R_1\leftrightarrow R_2}
\left[\begin{array}{rrr|rrr}1&1&1&0&1&0\\0&1&1&1&0&0\\3&4&3&0&0&1\end{array}\right]\\
&\xrightarrow{R_3\leftarrow R_3-3R_1}
\left[\begin{array}{rrr|rrr}1&1&1&0&1&0\\0&1&1&1&0&0\\0&1&0&0&-3&1\end{array}\right]\\
&\xrightarrow{\substack{R_3\leftarrow R_3-R_2\\ R_3\leftarrow -R_3}}
\left[\begin{array}{rrr|rrr}1&1&1&0&1&0\\0&1&1&1&0&0\\0&0&1&1&3&-1\end{array}\right]\\
&\xrightarrow{\substack{R_1\leftarrow R_1-R_3\\R_2\leftarrow R_2-R_3}}
\left[\begin{array}{rrr|rrr}1&1&0&-1&-2&1\\0&1&0&0&-3&1\\0&0&1&1&3&-1\end{array}\right]\\
&\xrightarrow{R_1\leftarrow R_1-R_2}
\left[\begin{array}{rrr|rrr}1&0&0&-1&1&0\\0&1&0&0&-3&1\\0&0&1&1&3&-1\end{array}\right].
\end{align*}
Thus
\[
A^{-1}=B=\begin{pmatrix}-1&1&0\\0&-3&1\\1&3&-1\end{pmatrix}.
\]
\notetarget{note:02:inverse-checks}Check both sides using rows:
\[
AB=\begin{pmatrix}\operatorname{row}_2(B)+\operatorname{row}_3(B)\\
\operatorname{row}_1(B)+\operatorname{row}_2(B)+\operatorname{row}_3(B)\\
3\operatorname{row}_1(B)+4\operatorname{row}_2(B)+3\operatorname{row}_3(B)
\end{pmatrix}=I_3,
\]
\[
BA=\begin{pmatrix}-\operatorname{row}_1(A)+\operatorname{row}_2(A)\\
-3\operatorname{row}_2(A)+\operatorname{row}_3(A)\\
\operatorname{row}_1(A)+3\operatorname{row}_2(A)-\operatorname{row}_3(A)
\end{pmatrix}=I_3.
\]
The full solution of \(Ax=b\) is therefore
\(x=(-b_1+b_2,-3b_2+b_3,b_1+3b_2-b_3)^\top\).
\end{example}

\begin{proposition}[Row operations preserve column relations]\label{prop:02:relations}\notetarget{note:02:column-relations}
If \(H=EA\) with \(E\) invertible, then \(N(H)=N(A)\).
More explicitly, writing the columns as \(h_j=Ea_j\),
\[
\sum_jc_jh_j=0\quad\Longleftrightarrow\quad\sum_jc_ja_j=0.
\]
\end{proposition}
\begin{proof}
The left sum equals \(E\sum_jc_ja_j\). A zero sum on the right gives a
zero sum on the left; multiplying a zero sum on the left by \(E^{-1}\)
gives a zero sum on the right. Taking arbitrary coefficient columns
\(c\) proves the null-space equality.
\end{proof}
Row operations may change the columns, but preserve every coefficient relation.

\needspace{16\baselineskip}
\section*{2. Block multiplication and trace}

\subsection*{Partitioned products}
\notetarget{note:02:block-products}Partition the rows and columns into consecutive, nonempty groups. Write
\[
 A=[A_{ij}]_{r\times s}\in\R^{m\times n},\qquad
 B=[B_{jk}]_{s\times t}\in\R^{n\times p},
\]
where \(r\times s\) and \(s\times t\) count the \emph{blocks}, and
\[
 A_{ij}\in\R^{m_i\times n_j},\qquad
 B_{jk}\in\R^{n_j\times p_k},\qquad
 m=\sum_{i=1}^r m_i,\quad n=\sum_{j=1}^s n_j,\quad p=\sum_{k=1}^t p_k.
\]
Thus the column groups of \(A\) match the row groups of \(B\),
with the same respective sizes \(n_1,\ldots,n_s\).
\notetarget{note:02:finite-blocks}The product has \(r\) block rows and \(t\) block columns:
\[
 \boxed{AB=C=[C_{ik}]_{r\times t},\qquad
 C_{ik}=\sum_{j=1}^s A_{ij}B_{jk}\in\R^{m_i\times p_k}.}
\]
All summands have the displayed size; the blocks need not be square.

\begin{proof}
\notetarget{note:02:block-proof}\begin{itemize}
\item Fix \(i,k\) and a within-block position
\(1\leqslant\alpha\leqslant m_i\), \(1\leqslant\beta\leqslant p_k\).
Let \(x\) be the corresponding full row of \(A\), and \(y\) the
corresponding full column of \(B\). Put
\[
 N_0=0,\qquad N_j=n_1+\cdots+n_j\quad(1\leqslant j\leqslant s).
\]
\item In the \(j\)th pair of matching groups,
\[
 x_{N_{j-1}+\ell}=(A_{ij})_{\alpha\ell},\qquad
 y_{N_{j-1}+\ell}=(B_{jk})_{\ell\beta}
 \quad(1\leqslant\ell\leqslant n_j).
\]
\item Group the defining scalar product by those same intervals:
\[
\begin{aligned}
 (C_{ik})_{\alpha\beta}
 &=\sum_{h=1}^{n}x_hy_h
 =\sum_{j=1}^{s}\sum_{\ell=1}^{n_j}
     x_{N_{j-1}+\ell}y_{N_{j-1}+\ell}\\
 &=\sum_{j=1}^{s}\sum_{\ell=1}^{n_j}
     (A_{ij})_{\alpha\ell}(B_{jk})_{\ell\beta}
 =\left(\sum_{j=1}^{s}A_{ij}B_{jk}\right)_{\alpha\beta}.
\end{aligned}
\]
Every entry agrees, proving the formula.
\end{itemize}
\end{proof}

\notetarget{note:02:two-blocks}For \(r=s=t=2\), the formula specializes to
\[
A=\begin{pmatrix}A_{11}&A_{12}\\A_{21}&A_{22}\end{pmatrix},\qquad
B=\begin{pmatrix}B_{11}&B_{12}\\B_{21}&B_{22}\end{pmatrix},
\]
\[
AB=\begin{pmatrix}
A_{11}B_{11}+A_{12}B_{21}&A_{11}B_{12}+A_{12}B_{22}\\
A_{21}B_{11}+A_{22}B_{21}&A_{21}B_{12}+A_{22}B_{22}
\end{pmatrix}.
\]

\notetarget{note:02:outer-products}In particular, if \(A=[\,a_1\ \cdots\ a_n\,]\) and the rows of \(B\)
are \(b_1^\top,\ldots,b_n^\top\), then
\[
AB=\sum_{j=1}^n a_jb_j^\top.
\]
The \((i,k)\) entry on the right is \(\sum_j a_{ij}b_{jk}\), exactly
that of \(AB\). Also \([A\ B]\binom{x}{y}=Ax+By\), by splitting the
column combination into the two groups.

\begin{example}[A block system]\label{example:02:block}\notetarget{note:02:block-example}
Let \(A\in\R^{r\times r}\), \(B\in\R^{r\times s}\),
\(D\in\R^{s\times s}\), and suppose \(A,D\) are invertible.
Block multiplication writes
\[
\begin{pmatrix}A&B\\0&D\end{pmatrix}\binom{x}{y}=\binom{u}{v}
\quad\Longleftrightarrow\quad
Ax+By=u,\quad Dy=v.
\]
First \(y=D^{-1}v\), then \(x=A^{-1}u-A^{-1}BD^{-1}v\).
Thus the candidate inverse is
\[
G=\begin{pmatrix}A^{-1}&-A^{-1}BD^{-1}\\0&D^{-1}\end{pmatrix}.
\]
\notetarget{note:02:block-checks}Both orders give the identity:
\[
\begin{pmatrix}A&B\\0&D\end{pmatrix}G
=\begin{pmatrix}I_r&-BD^{-1}+BD^{-1}\\0&I_s\end{pmatrix}=I_{r+s},
\]
\[
G\begin{pmatrix}A&B\\0&D\end{pmatrix}
=\begin{pmatrix}I_r&A^{-1}B-A^{-1}B\\0&I_s\end{pmatrix}=I_{r+s}.
\]
Hence \(G\) is the inverse.
\end{example}

\subsection*{Trace}
\notetarget{note:02:trace}For \(A\in\R^{n\times n}\), its \textbf{trace} is the sum of its
\textbf{diagonal entries}:
\[
\tr A=\sum_{i=1}^n a_{ii}.
\]
A matrix is \textbf{diagonal} when all entries off its main diagonal
are zero; \(\operatorname{diag}(d_1,\ldots,d_n)\) denotes the matrix
with those diagonal entries.

\begin{proposition}[Trace identities]\label{prop:02:trace}
For square matrices of the same size,
\[
\tr(sA+tB)=s\tr A+t\tr B,\qquad \tr(A^\top)=\tr A.
\]
\notetarget{note:02:trace-cyclic}For \(A\in\R^{m\times n}\), \(B\in\R^{n\times m}\),
\[
\tr(AB)=\tr(BA).
\]
\end{proposition}
\begin{proof}
The first identity is \(\sum_i(sa_{ii}+tb_{ii})=s\sum_i a_{ii}+t\sum_i b_{ii}\).
Transposition leaves every diagonal entry unchanged, proving the second.
For the last identity, the two square products may have different sizes,
but
\[
\tr(AB)=\sum_{i=1}^m\sum_{j=1}^n a_{ij}b_{ji}
=\sum_{j=1}^n\sum_{i=1}^m b_{ji}a_{ij}=\tr(BA).
\]
\end{proof}
If \(A,B,C\) have sizes \(m\times n\), \(n\times p\), \(p\times m\),
apply the two-factor result to \(A\) and \(BC\), then to \(B\) and \(CA\),
to obtain
\[
\tr(ABC)=\tr(BCA)=\tr(CAB).
\]

\begin{proposition}[A trace that detects zero]\label{prop:02:squares}\notetarget{note:02:trace-squares}
For a real \(m\times n\) matrix \(A\),
\[
\tr(A^\top A)=\sum_{i=1}^m\sum_{j=1}^n a_{ij}^2,
\qquad \tr(A^\top A)=0\quad\Longleftrightarrow\quad A=0.
\]
\end{proposition}
\begin{proof}
The \(j\)th diagonal entry of \(A^\top A\) is \(\sum_i a_{ij}^2\);
summing over \(j\) proves the formula. Every summand is nonnegative.
Their sum is zero exactly when each entry \(a_{ij}\) is zero.
\end{proof}

\begin{example}[An impossible matrix equation]\label{example:02:commutator}\notetarget{note:02:commutator}
Can real \(3\times3\) matrices \(X,Y\) satisfy
\[
XY-YX=\operatorname{diag}(1,0,0)?
\]
The left side has trace \(\tr(XY)-\tr(YX)=0\); the right side has
trace \(1\). Equality would force \(0=1\), so no such matrices exist.
\end{example}

\section*{3. Linear spaces and subspaces}

\notetarget{note:02:linear-space}A \textbf{real linear space} of matrices is a nonempty collection of
matrices of one fixed size, closed under addition and multiplication by
real numbers. Its elements are also called vectors. The addition and
scalar multiplication are the usual entrywise operations.
Such a collection contains zero: choose one member \(X\), then
\(0X=0\). It contains \(-X=(-1)X\) whenever it contains \(X\).
The eight coordinate laws from Unit 1 still hold; closure keeps every
sum and scalar multiple in the collection.

The full collection \(\R^{m\times n}\) is a linear space: sums and
scalar multiples have the same size. This includes the spaces of columns \(\R^m\) and of scalars \(\R\). \notetarget{note:02:polynomials}For real polynomials of degree
at most \(d\), including the zero polynomial, use the coefficient column
\[
p(t)=a_0+a_1t+\cdots+a_dt^d
\quad\longleftrightarrow\quad(a_0,\ldots,a_d)^\top.
\]
Here a polynomial means its finite coefficient expression; equality
means equality of coefficients. Addition and scaling are coordinatewise,
so this is another notation for a linear space, denoted by \(\mathcal P_d\).

\notetarget{note:02:subspace-test}A \textbf{subspace} of a linear space \(V\) is a subset that is itself
a linear space with the inherited operations.

\begin{proposition}[Subspace test]\label{prop:02:subspace}
A subset \(W\subseteq V\) is a subspace if and only if it is nonempty
and
\[
u,v\in W,\quad s,t\in\R\quad\Longrightarrow\quad su+tv\in W.
\]
\end{proposition}
\begin{proof}
A subspace is nonempty and is closed under each scalar multiplication
and then addition, so it satisfies the displayed condition.
Conversely choose \(u\in W\). Taking both coefficients zero gives
\(0\in W\). Taking \(s=t=1\) proves addition closure; taking \(t=0\)
proves scalar closure. The coordinate laws are inherited from \(V\),
so \(W\) is a linear space.
\end{proof}
\notetarget{note:02:zero-whole-subspaces}Both \(\{0\}\) and \(V\) satisfy this test. A subspace must contain zero,
but this alone is insufficient: \(\{0,1\}\subseteq\R\) contains zero
and fails addition closure because \(1+1=2\notin\{0,1\}\).

\begin{example}[Homogeneous conditions]\label{example:02:homogeneous-space}\notetarget{note:02:homogeneous-example}
In \(\R^4\), consider
\[
W=\{(a,b,c,d)^\top:a+2b=c-d\}.
\]
Define \(f(a,b,c,d)=a+2b-c+d\). Entrywise addition and scaling give
\[
f(su+tv)=sf(u)+tf(v).
\]
Thus \(f(0)=0\), and \(f(u)=f(v)=0\) implies \(f(su+tv)=0\).
The subspace test proves that \(W\) is a subspace. \notetarget{note:02:homogeneous-description}Solving for \(a\) gives its complete description
\[
W=\left\{
 b\begin{pmatrix}-2\\1\\0\\0\end{pmatrix}
+c\begin{pmatrix}1\\0\\1\\0\end{pmatrix}
+d\begin{pmatrix}-1\\0\\0\\1\end{pmatrix}:b,c,d\in\R
\right\}.
\]
Every vector in the displayed set satisfies the condition, and every
member of \(W\) is obtained using its own last three coordinates.

\notetarget{note:02:subspace-examples}In contrast, the set of \(n\times n\) matrices satisfying \(A^2=I_n\)
for \(n\geqslant1\) is not a subspace: \(0^2\ne I_n\).
\end{example}

Unit 1, Proposition~\ref{one-prop:01:closure}, proves that
\(\mathcal C(A)\subseteq\R^m\) and \(N(A)\subseteq\R^n\) are subspaces.
\notetarget{note:02:symmetric-subspace}The set of symmetric \(n\times n\) matrices is another example:
\(0^\top=0\), and for symmetric \(S,T\),
\((sS+tT)^\top=sS^\top+tT^\top=sS+tT\).

\notetarget{note:02:homogeneous-matrix-space}For fixed \(A\in\R^{m\times n}\), \(B\in\R^{p\times q}\),
the set \(\{X\in\R^{n\times p}:AXB=0\}\) is a subspace of
\(\R^{n\times p}\): it contains zero, and
\[
A(sX+tY)B=sAXB+tAYB=0
\]
whenever \(AXB=AYB=0\).

\subsection*{Spanning sets}
\notetarget{note:02:span}For a finite list \(v_1,\ldots,v_k\in V\), define its \textbf{span} by
\[
\operatorname{span}(v_1,\ldots,v_k)
=\left\{\sum_{j=1}^k c_jv_j:c_1,\ldots,c_k\in\R\right\}.
\]
The span of the empty list is \(\{0\}\). A list \textbf{spans} \(W\)
when this set equals \(W\); its members are then \textbf{generators}
of \(W\). Thus \(\mathcal C(A)\) is the span of the columns of \(A\).

\begin{proposition}[The smallest subspace containing a list]\label{prop:02:span}
The span of \(v_1,\ldots,v_k\) is a subspace containing every \(v_j\).
It is contained in every subspace containing the list.
\end{proposition}
\begin{proof}
The all-zero coefficients give zero. If \(x=\sum_j a_jv_j\) and
\(y=\sum_j b_jv_j\), then
\(sx+ty=\sum_j(sa_j+tb_j)v_j\), proving closure.
\notetarget{note:02:span-minimality}Putting coefficient \(1\) on \(v_j\) and \(0\) on the others gives
\(v_j\) itself. Finally a subspace containing all the \(v_j\) contains
all their finite linear combinations, by repeated addition and scaling.
For an empty list its span is \(\{0\}\), which is in every subspace.
\end{proof}

\begin{proposition}[Removing a redundant generator]\label{prop:02:redundancy}\notetarget{note:02:redundancy}
If \(v_j\) is a combination of the other members of a generating list,
removing it does not change the span.
\end{proposition}
\begin{proof}
Write \(v_j=\sum_{i\ne j}d_iv_i\). Every combination of the full list is
\[
\sum_i c_iv_i=\sum_{i\ne j}(c_i+c_jd_i)v_i,
\]
so it is in the shorter span. Conversely a combination of the shorter
list is a combination of the original list with coefficient zero on
\(v_j\). These give the two inclusions.
\end{proof}

\begin{example}[Matrices commuting with a fixed matrix]\label{example:02:commuting}\notetarget{note:02:commuting-equations}
Find all matrices \(B\) commuting with
\(A=\begin{pmatrix}1&2\\3&4\end{pmatrix}\).
Writing \(B=\begin{pmatrix}a&b\\c&d\end{pmatrix}\), multiplication gives
\[
AB=\begin{pmatrix}a+2c&b+2d\\3a+4c&3b+4d\end{pmatrix},\qquad
BA=\begin{pmatrix}a+3b&2a+4b\\c+3d&2c+4d\end{pmatrix}.
\]
Equality is equivalent to
\[
2c=3b,\qquad 2d=2a+3b,\qquad c=d-a.
\]
\notetarget{note:02:commuting-solutions}The first two give \(c=3b/2\), \(d=a+3b/2\), which also satisfy the
third. Set \(\alpha=a\), \(\beta=b/2\). All solutions are exactly
\[
B=\begin{pmatrix}\alpha&2\beta\\3\beta&\alpha+3\beta\end{pmatrix}
=\alpha I_2+\beta\begin{pmatrix}0&2\\3&3\end{pmatrix}.
\]
Conversely these entries satisfy all three equations, hence \(AB=BA\).
\notetarget{note:02:commuting-closure}The solution set is the span of the two displayed matrices and is
therefore a subspace. A direct check also works: \(A0=0A\), and if \(AB=BA\), \(AC=CA\), then
\[
A(sB+tC)=sAB+tAC=sBA+tCA=(sB+tC)A.
\]
\end{example}

\section*{4. Intersections, sums, and translates}

\begin{proposition}[Intersection and sum]\label{prop:02:sum-intersection}\notetarget{note:02:sum-intersection}
If \(U,W\) are subspaces of the same \(V\), then so are
\[
U\cap W,\qquad U+W:=\{u+w:u\in U,\ w\in W\}.
\]
Moreover, \(U+W\) is the smallest subspace containing both \(U\) and \(W\).
\end{proposition}
\begin{proof}
Zero belongs to both sets. If \(x,y\in U\cap W\), then \(sx+ty\)
belongs to \(U\) and to \(W\), hence to the intersection.
\notetarget{note:02:sum-proof}For \(x=u_1+w_1\), \(y=u_2+w_2\),
\[
sx+ty=(su_1+tu_2)+(sw_1+tw_2)\in U+W.
\]
Thus the sum is also a subspace. \notetarget{note:02:sum-minimality}Taking \(w=0\) or \(u=0\) shows that
it contains \(U\) and \(W\). Any subspace containing both must contain
all their sums, giving the asserted minimality.
\end{proof}

\begin{proposition}[When is a union a subspace?]\label{prop:02:union}\notetarget{note:02:union}
For subspaces \(U,W\) of \(V\), the union \(U\cup W\) is a subspace
if and only if \(U\subseteq W\) or \(W\subseteq U\).
\end{proposition}
\begin{proof}
\begin{itemize}
\item If either inclusion holds, the union is the larger subspace.
\item \notetarget{note:02:union-necessity}If neither holds, choose \(u\in U\setminus W\), \(w\in W\setminus U\).
Then
\[
u+w\in U\ \Longrightarrow\ w=(u+w)-u\in U,\qquad
u+w\in W\ \Longrightarrow\ u=(u+w)-w\in W.
\]
Both contradict the choices. Hence \(u,w\in U\cup W\), but
\(u+w\notin U\cup W\): addition closure fails.
\end{itemize}
\end{proof}

\begin{proposition}[Uniqueness of a sum decomposition]\label{prop:02:direct-sum}\notetarget{note:02:direct-sum}
Let \(U,W\) be subspaces of the same real linear space \(V\).
Every \(x\in U+W\) has a unique representation \(x=u+w\) with
\(u\in U\), \(w\in W\), if and only if \(U\cap W=\{0\}\).
When this holds the sum is called a \textbf{direct sum}, written \(U\oplus W\).
\end{proposition}
\begin{proof}
Existence follows from the definition of \(U+W\).
If \(u+w=u'+w'\), then \(u-u'=w'-w\in U\cap W\).
A zero intersection forces \(u=u'\), \(w=w'\).
\notetarget{note:02:direct-sum-converse}Conversely, for \(z\in U\cap W\), the two representations
\(z=z+0=0+z\) must agree term by term. Hence \(z=0\).
\end{proof}

\begin{example}[Symmetric and skew-symmetric parts]\label{example:02:decomposition}\notetarget{note:02:symmetric-skew-spaces}
A matrix \(K\) is \textbf{skew-symmetric} if \(K^\top=-K\).
The symmetric and skew-symmetric \(n\times n\) matrices each form a
subspace: transpose linearity proves closure, and zero belongs to both.
\notetarget{note:02:symmetric-skew-construction}Every \(A\in\R^{n\times n}\) can be written
\[
A=S+K,\qquad S=\frac{A+A^\top}{2},\qquad K=\frac{A-A^\top}{2}.
\]
Indeed their sum is \(A\), while the transpose laws give
\(S^\top=S\), \(K^\top=-K\). \notetarget{note:02:symmetric-skew-uniqueness}A matrix in both subspaces satisfies
\(X=X^\top=-X\), so \(2X=0\) and every entry is zero.
Proposition~\ref{prop:02:direct-sum} therefore proves that this
decomposition is unique.
\end{example}

\notetarget{note:02:translate}For \(p\in V\) and a subspace \(W\), the set
\(p+W=\{p+w:w\in W\}\) is a \textbf{translate} of \(W\), also called
an \textbf{affine subspace}.

\begin{proposition}[A translate that is a subspace]\label{prop:02:translate}
The set \(p+W\) is a subspace if and only if \(p\in W\); in that case
\(p+W=W\).
\end{proposition}
\begin{proof}
If \(p\in W\), then \(p+w\in W\) for all \(w\in W\). Conversely every
\(w\in W\) equals \(p+(w-p)\), with \(w-p\in W\), proving equality.
If \(p+W\) is a subspace, it contains zero, so \(p+w=0\) for some
\(w\in W\). Hence \(p=-w\in W\).
\end{proof}
\notetarget{note:02:affine-system}For a consistent system \(Ax=b\), Unit 1,
Theorem~\ref{one-thm:01:translate}, gives the affine subspace
\(x_0+N(A)\). It is a linear subspace exactly when \(b=0\): if it is a
subspace then zero solves the equation, forcing \(b=A0=0\); if \(b=0\)
its solution set is the subspace \(N(A)\).

\notetarget{note:02:affine-matrix-system}For fixed \(A\in\R^{m\times n}\), \(B\in\R^{p\times q}\),
\(C\in\R^{m\times q}\), suppose \(X_0\in\R^{n\times p}\) satisfies
\(AX_0B=C\). Then, with \(X,Z\in\R^{n\times p}\),
\[
\{X:AXB=C\}=X_0+\{Z:AZB=0\}.
\]
Indeed \(AXB=C\) implies \(A(X-X_0)B=0\), and conversely
\(AZB=0\) implies \(A(X_0+Z)B=C\).
